ここでは高被引用論文の変遷をさらに詳しく見るために論文順位の変動を観察する。
Kendallの順位相関を使いた比較を年ごとに行えば、相関係数が高いほど順位変動が少ない、すなわち安定していると見做せる。
対象はすべての累積世代の高被引用論文（536論文）とし、1947-2020年までの各年につき、前年と比較した際の相関係数を算出した。
図\ref{fig:gr101}に各年の相関係数を示す。
値は0.8-0.95の間を推移し、経過年による傾向は見られない。
さらに、累積的な被引用の効果がどのように現れるかを確認するため、出版後n年までの累積被引用数を用いた場合と比較した。
完全な累積型との相関係数の差を図\ref{fig:gr102}に示す。
これによれば、影響年を制限するほど相関係数が高くなる、すなわち順位が安定する傾向があり、被引用における累積的な影響は順位を変動させる要因となる可能性が高い。
